\section{Модуль 2}

% Дані розділу — Semester_Scope.md (обсяг робіт на модуль 2, автор — П5).

% ===== Ціль MVP і що переходить у бету =====
% Semester_Scope.md, розділи 1, 3 і 4.
\begin{frame}{Обсяг робіт на модуль 2}
\vspace{0.05cm}
\lbl{Ціль MVP}\par\vskip0.15cm
\begin{tikzpicture}
  \node[anchor=north west,inner sep=0pt,text width=13.9cm,align=left] (g) at (0.3,0)
    {{\footnotesize\semibold\color{ink}\linespread{1.2}\selectfont Пілот завантажує відео GoPro з «Жаги швидкості» і в~браузері бачить свої кола, повороти, де втратив найбільше часу, і~2–3 поради з~цифрами; час кола відрізняється від офіційного хронометражу в~середньому не більше ніж на 0,05\,с.\par}};
  \fill[kerb] (0,0) rectangle ($(g.south west)+(-0.24,0)$);
\end{tikzpicture}

\vspace{0.2cm}
\begin{columns}[T,onlytextwidth]
\column{0.17\textwidth}\stat{18}{задач у MVP на \textbf{48 год} — рівно ресурс}
\column{0.17\textwidth}\stat[ink]{−8}{год: 4 спрощення Must — \id{B-05}, \id{B-08}, \id{B-09}, \id{B-13}}
\column{0.17\textwidth}\stat[lane]{15}{год: 6 задач переходять у бету}
\column{0.43\textwidth}
  \lbl[lane]{Переходить у бету}\par\vskip0.1cm
  \renewcommand{\arraystretch}{1.12}
  {\tiny
  \begin{tabularx}{\linewidth}{@{}>{\ttfamily}lXr@{}}
  B-08 & Карта траси з траєкторією (Must) & 3 \\
  B-03 & Повідомлення для камер без GPS & 2 \\
  B-06 & Межі траси за супутником & 4 \\
  B-12 & Відкидання кіл із трафіком & 2 \\
  B-16 & Мітка сили зв'язку біля поради & 2 \\
  B-24 & Перевірка Tenth Finder і Kart Track & 2 \\
  \end{tabularx}}
  \vskip0.05cm
  {\tiny\color{mute}Окрім \id{B-08}, усі — Should: вони переходять першими.\par}
\end{columns}
\note{\speaker{П5}
Ціль MVP одним реченням: пілот завантажує відео GoPro з «Жаги швидкості» і в браузері бачить свої кола, повороти, де втратив найбільше часу, і 2–3 поради з цифрами. Час кола має відрізнятися від офіційного хронометражу в середньому не більше ніж на 0,05 секунди.

У MVP — 18 задач на 48 годин, рівно наш ресурс. Щоб вкластися, ми спростили чотири обов'язкові задачі й зекономили 8 годин. Шість задач на 15 годин переходять у бету: це п'ять задач Should і карта траси.}
\end{frame}

% ===== Розподіл по людях =====
% Semester_Scope.md, розділ 5. Червоним — задачі критичного шляху (розділ 5, «Порядок»).
\begin{frame}{Розподіл по людях: по 12 годин}
\vspace{0.2cm}
\centering
\begin{tikzpicture}
  \node[anchor=west,font=\tiny\semibold,text=mute] at (3.3,0.62) {\caps{Задачі MVP, ширина — оцінка в годинах}};
  \node[anchor=east,font=\tiny\semibold,text=mute] at (14.2,0.62) {\caps{Разом, год}};
  \pplrow{1}{У1}{Платформа і завантаження}{B-18/3/ink,B-01/4/ink,B-02/2/ink}{3}
  \pplrow{2}{У2 · лід}{Дані, точність, демо}{B-19/2/ink,B-20/3/kerb,B-21/2/ink,B-23/1/ink,B-22/2/kerb}{2}
  \pplrow{3}{У3}{Кола і відрізки}{B-04/4/kerb,B-05/1/ink,B-10/4/kerb}{3}
  \pplrow{4}{У4}{Повороти і втрати}{B-07/2/ink,B-09/4/kerb,B-11/3/ink,B-13/1/ink}{2}
  \pplrow{5}{У5}{Поради й екран розбору}{B-14/5/kerb,B-15/3/kerb,B-17/2/kerb}{2}
  \draw[line] (0,-3.88) -- (14.2,-3.88);
  \node[anchor=west,font=\footnotesize\semibold,text=ink] at (0,-4.18) {18 задач};
  \node[anchor=west,font=\tiny,text=body] at (3.3,-4.18) {%
    \tikz[baseline=-0.6ex]\fill[kerb] (0,-0.1) rectangle (0.22,0.12);\;критичний шлях, 27 год\qquad
    \tikz[baseline=-0.6ex]\fill[ink] (0,-0.1) rectangle (0.22,0.12);\;інші задачі\qquad
    \tikz[baseline=-0.6ex]\fill[soft] (0,-0.1) rectangle (0.22,0.12);\;запас};
  \node[anchor=east,font=\tiny,text=mute] at (14.2,-4.18) {48 на задачі + 12 запас = 60};
\end{tikzpicture}
\note{\speaker{П5}
У команді п'ять потоків, кожен — приблизно 12 годин за два тижні: 9–10 годин на задачі і 2–3 години запасу. Ширина блоку — оцінка задачі в годинах.

Перший учасник робить платформу і завантаження відео. Лід — дані, точність і демо: формат даних, тестові дані і фінальний прогін зводять роботу всіх. Третій — поділ на кола і відрізки, четвертий — повороти і втрати часу, п'ятий — поради й екран розбору.

Червоним позначено критичний шлях. У MVP він 27 годин, бо метрики повороту ми спростили на годину.}
\end{frame}

% ===== Порядок робіт =====
% Semester_Scope.md, розділ 5 «Порядок». Відкриті питання — Project-Overview.md і звірка артефактів.
\begin{frame}{Порядок на два тижні}
\vspace{0cm}
\centering
\begin{tikzpicture}
  \def\dw{0.8875}% см на день, 4–20 листопада
  % тижні
  \fill[ink] (0,0) rectangle ({7*\dw-0.04},0.42);
  \fill[kerb] ({8*\dw},0) rectangle ({16*\dw},0.42);
  \node[anchor=west,font=\tiny\semibold,text=white] at (0.1,0.21) {ТИЖДЕНЬ 1\enspace·\enspace 4–11 листопада};
  \node[anchor=west,font=\tiny\semibold,text=white] at ({8*\dw+0.1},0.21) {ТИЖДЕНЬ 2\enspace·\enspace 12–20 листопада};
  % перші 2 дні
  \draw[mute,line width=0.5pt] (0,0.55) -- (0,0.75) -- ({2*\dw},0.75) -- ({2*\dw},0.55);
  \node[anchor=south west,font=\tiny,text=body,inner sep=1pt] at (0,0.78) {перші 2 дні: \id{B-18}, \id{B-19}, \id{B-20}};
  % дедлайн прогону
  \draw[kerb,dashed,line width=0.7pt] ({14*\dw},-0.1) -- ({14*\dw},0.95);
  \node[anchor=south east,font=\tiny,text=mute,inner sep=1pt] at ({16*\dw},0.78) {{\semibold\color{kerb}\id{B-22} до 18.11}\enspace·\enspace 2 дні на виправлення};
  % що в тижнях
  \node[anchor=north west,font=\tiny,text=body,text width=6.0cm,align=left,inner sep=0pt] at (0,-0.15)
    {\id{B-23} — у перший тиждень, як вимагає Project Overview.\\ Паралельно: \id{B-01}, \id{B-02}; \id{B-04}, \id{B-05}; \id{B-07}; \id{B-21} після \id{B-04}.};
  \node[anchor=north west,font=\tiny,text=body,text width=6.9cm,align=left,inner sep=0pt] at ({8*\dw},-0.15)
    {\id{B-10} → \id{B-09}, \id{B-11}, \id{B-13} → \id{B-14}, \id{B-15}, \id{B-17} → \id{B-22}};
\end{tikzpicture}

\vspace{0.2cm}
\raggedright
\begin{columns}[T,onlytextwidth]
\column{0.47\textwidth}
  \lbl{Критичний шлях: 27 год}\par\vskip0.12cm
  {\tiny\color{body}При 6 год на тиждень на людину він не влазить у 2 тижні, якщо чекати готового коду попередника.\par\vskip0.1cm
  Тому кожен модуль пишемо і тестуємо на тестових даних демо (\id{B-20}) у форматі \id{B-19}. У4 і У5 стартують у першому тижні з вручну підготовленого JSON сесії. Зводимо все в \id{B-22}.\par}
\column{0.47\textwidth}
  \lbl[ink]{Відкриті питання}\par\vskip0.12cm
  {\tiny\color{body}
  \begin{itemize}\setlength{\itemsep}{0.06cm}
    \item \textbf{Технології не обрані.} Поки припускаємо TypeScript у браузері, gopro-telemetry і MVP без сервера.
    \item \textbf{Tenth Finder і Kart Track.} Overview перевіряє їх у перший тиждень, а в обсязі MVP \id{B-24} — у беті.
    \item \textbf{Тривалість.} 4–20 листопада — це 2,4 тижня, а не 2. Якщо рахувати весь період, на задачі буде 58 год.
  \end{itemize}}
\end{columns}
\note{\speaker{Лід · Яна}
Перші два дні — каркас, формат даних і перевірка gopro-telemetry: від них залежать усі. Прогін демо — до 18 листопада, щоб лишилося два дні на виправлення.

Критичний шлях — 27 годин: якщо чекати код попередника, він не влазить. Тому кожен тестує свій модуль на даних демо у спільному форматі, а зводимо все в кінці.

Відкриті питання — на слайді: технології, перевірка конкурентів і тривалість модуля.}
\end{frame}

% ===== Definition of Done =====
% Semester_Scope.md, розділ 6.
\begin{frame}{Definition of Done}
\vspace{0.2cm}
\begin{columns}[T,onlytextwidth]
\column{0.47\textwidth}
  \lbl[ink]{Задача готова, коли}\par\vskip0.25cm
  {\scriptsize\color{body}
  \dod{Код у головній гілці, пройшов рев'ю одного учасника команди.}
  \dod{Є автотести на тестових даних демо (\id{B-20}); CI зелений.}
  \dod{Вхід і вихід модуля відповідають формату \id{B-19}.}
  \dod{Тексти для пілота — українською, за tone of voice з Product Map.}
  \dod{Відео не покидає комп'ютер пілота: на мережу йде лише код застосунку.}
  \dod{Задеплоєно на хостинг (\id{B-18}), якщо задача видна користувачу.}}
\column{0.47\textwidth}
  \lbl{MVP готовий, коли}\par\vskip0.25cm
  {\scriptsize\color{body}
  \dod{Усі 18 задач MVP відповідають DoD задачі.}
  \dod{Автотест \id{B-21} проходить: похибка часу кола проти AMB у середньому ≤\,0,05\,с на демо-сесії.}
  \dod{Сценарій демо проходить від початку до кінця на задеплоєній версії в Chrome на ноутбуці без помилок у консолі.}
  \dod{Задачі, що переїжджають у бету, внесені в беклог бети.}}
\end{columns}
\note{\speaker{П5}
Definition of Done — коли ми вважаємо роботу завершеною. Задача готова, коли код у головній гілці після рев'ю, є автотести на даних демо і CI зелений, модуль працює у спільному форматі даних, тексти українською, відео не покидає комп'ютер пілота, а все видиме користувачу задеплоєно.

MVP готовий, коли всі 18 задач виконані за цим правилом, автотест точності проходить, сценарій демо проходить на задеплоєній версії без помилок, а все, що переїжджає, внесено в беклог бети.}
\end{frame}

% ===== Сценарій демо =====
% Semester_Scope.md, розділ 7.
\begin{frame}{Сценарій демо: \textasciitilde5 хвилин}
\vspace{-0.05cm}
{\scriptsize\color{mute}Сесія «Жага швидкості», 1 жовтня 2026\enspace·\enspace GoPro HERO7, \id{GH011616} і \id{GH021616}\enspace·\enspace 14 кіл\par}
\vspace{0.2cm}
\centering
\begin{tikzpicture}
  \demostep{0}{0}{1}{Відкриваємо застосунок за посиланням}{Екран завантаження.}
  \demostep{3.6}{0}{2}{Обираємо обидва файли, показуємо вкладку мережі}{Одна сесія з двох частин; відео не вивантажується.}
  \demostep{7.2}{0}{3}{Таблиця кіл}{14 залікових кіл, виїзне відкинуто; найкраще — коло 11, 41,474\,с.}
  \demostep{10.8}{0}{4}{Результат автотесту \id{B-21}}{Похибка проти AMB у середньому ≤\,0,05\,с (у демо — 0,035\,с).}
  \demostep{0}{-2.12}{5}{Список поворотів T1–T7}{Відсортовано за втраченим часом, найбільші втрати першими.}
  \demostep{3.6}{-2.12}{6}{Картки порад}{2–3 поради з цифрою «зараз» і ціллю.}
  \demostep{7.2}{-2.12}{7}{Клікаємо пораду}{Перехід до цього повороту у списку.}
  \demostep{10.8}{-2.12}{8}{Що далі в беті}{Карта траси, відео біля поради, трек із телефона, оцінка порад.}
\end{tikzpicture}

\vspace{0.15cm}
\raggedright
{\tiny\color{body}\textbf{\color{kerb}Запасний план:} якщо завантаження не працює на демо-ноутбуці — відкриваємо збережений JSON сесії (\id{B-20}) і показуємо кроки 3–7.\par}
\note{\speaker{П5}
Демо MVP — приблизно п'ять хвилин. Беремо реальну сесію з «Жаги швидкості» — два файли GoPro, 14 кіл.

Обираємо обидва файли і показуємо, що відео не вивантажується. Далі — таблиця кіл, автотест точності, повороти за втраченим часом і картки порад; клік на пораду веде до повороту. Наприкінці — що буде в беті.

Якщо не спрацює завантаження, відкриваємо збережений JSON сесії.}
\end{frame}
